\section{Discussion}
\label{sec:discussion}

\subsection{Interpretation of the Component-wise Observations}

The ablation of Section V-C establishes that disabling any one of the three components degrades all three boundary measures simultaneously, while the region measures remain within a narrow band and are not maximised by the complete model. These are marginal effects measured within the Proposed model, each ablated configuration differing from it by a single disabled component; they characterise the role of each component within that model, and are not comparisons with the other members of the MAFF-Net family. The experiments establish that this association exists but do not, by themselves, identify its mechanism, since no analysis of the intermediate representations was undertaken. The accounts below are therefore interpretations consistent with the design of Section III and the observed behaviour, advanced as such rather than demonstrated explanations.

The Frequency Branch applies a fixed radial high-pass filter to each encoder feature and returns the retained content to the spatial domain. Intensity transitions at the lesion margin lie in the higher-frequency spectrum, and a plausible reading of the ablation is that an explicit route for this content preserves margin-associated information that a single representation would otherwise carry with the lower-frequency content of the uniform interior. One aspect of this comparison is unusually clean: the module has no trainable parameters, so replacing it with an identity mapping leaves the network's parameter count unchanged. The degradation that follows its removal therefore cannot be attributed to reduced capacity; in so far as this single-run difference reflects a genuine effect rather than run-to-run variation, it is attributable to the spectral operation itself. This does not establish which spectral content is responsible, only that the operation, not the capacity, would account for such a difference.

Removing the Feature Alignment Module produced the largest degradation in the average symmetric surface distance and the Boundary IoU of the three ablations. The two streams reaching the fusion stage differ in provenance: the spatial stream carries the full encoder feature, whereas the frequency stream, its low-frequency components suppressed, is approximately zero-mean and occupies a different range of values. A reasonable interpretation is that combining representations of dissimilar offset and scale lets one dominate, and that normalising both before the gate is computed allows the weighting to reflect the streams' content rather than their relative magnitudes. This carries an important qualification: the module comprises a normalisation, a learned projection, and a nonlinearity, and disabling it removes all three, so the ablated configuration has slightly fewer trainable parameters than the complete model. The observed degradation cannot be separated, on the present evidence, into the contribution of the normalisation and that of the reduced capacity.

Replacing the Gated Fusion Module with a fixed unweighted mean produced the largest increase in the 95th-percentile Hausdorff distance of the three ablations, but the smallest changes in the average symmetric surface distance and the Boundary IoU. A fixed mean applies the same weighting at every position and input, whereas the gate is computed from the streams themselves and may vary across channels and positions; if the relative usefulness of contextual and margin-associated content differs between the interior and the margin, a position-dependent weighting can reflect that where a fixed rule cannot. That its effect is most pronounced on the 95th-percentile distance, governed by the larger contour deviations, and least on the mean surface distance may indicate that its influence is concentrated on the more severe departures from the reference contour --- a conjecture the present experiments do not test. Here too the ablated configuration has marginally fewer trainable parameters.

\subsection{Boundary Fidelity versus Region Overlap}

That the improvements associated with these components appear in the boundary measures and not in the Dice coefficient follows from a geometric property of the measures, not from any peculiarity of the models: region overlap is area-based, so displacing a contour by a few pixels alters the overlapping area --- and the Dice coefficient --- only slightly, whereas the same displacement enters a contour-distance measure directly and in full. A modification acting principally on the placement of the margin may therefore leave region overlap almost unchanged while producing a measurable change in the boundary measures. The observations of Section V-C are consistent with this: the region measures varied within a narrow band across the four configurations, and the complete model recorded the highest value on none of them, whereas the boundary measures separated the configurations consistently.

Two consequences follow. First, the contribution of the three components is properly characterised as an increase in the fidelity of the predicted contour relative to the configurations in which they are disabled, not as an increase in region-overlap performance; the evidence does not support the latter and no such claim is made. This characterisation is internal to the Proposed model, and is not a claim of superior contour fidelity relative to the other members of the MAFF-Net family, among which, as recorded in Section VI-D, the Restored MAFF-Net attains the closer agreement with the reference contours. Second, region-overlap measures alone are insufficient to characterise methods designed to act on the lesion margin, a point for which Section V-B provides independent support in the disagreement between the two families' orderings. This observation is methodological rather than architectural: it follows from the geometric property noted above rather than any feature of the models, and is therefore expected to bear on the evaluation of other methods intended to act on the object boundary. Both consequences should be weighed against the magnitudes involved, which are small relative to the per-image dispersion of the distance measures.

\subsection{Retaining the Frozen Foundation Backbone}

Retaining the pretrained encoder as a fixed feature extractor carries several advantages here. Its representations are learned from a volume and diversity of data the available dermoscopic collections do not approach, which is material given the size of the datasets employed. Because the encoder is not updated, the trainable capacity is concentrated in the downstream modules and forms a small fraction of the total parameter count, reducing the number of parameters that must be estimated from the comparatively limited annotated data.

A further advantage is methodological. Because the same frozen encoder is used by the Proposed MAFF-Net, the Restored MAFF-Net, and every ablated configuration, the features entering the fusion pathway are identical across them, so differences in behaviour can be attributed to the modules under study rather than feature extraction, which permits the ablation of Section V-C to be read as an assessment of the components themselves; an architecture that also adapted the backbone would not support that inference as directly. The same reasoning applies to the Proposed-versus-Restored comparison: because the two share the frozen encoder and differ only in the fusion pathway, the difference between them in Sections V-A and V-B is attributable to the fusion design rather than feature extraction. On these benchmarks the dual-stream fusion does not surpass the multi-scale perception adaptation and frequency-guided feature fusion of the Original MAFF-Net, and the dual-stream reformulation is accordingly to be understood as an alternative fusion design examined and analysed here rather than one demonstrated to improve upon the original.

These advantages should not be overstated. The experiments include no configuration in which the encoder is fine-tuned, so the comparative merits of freezing and fine-tuning are not assessed here; the rationale rests on the considerations above rather than on evidence produced by this study. As noted in Section V-A, the Paper-faithful Reconstruction was trained without pretrained encoder weights, so the comparison between it and the remaining members of the family does not isolate the effect of the pretrained backbone.

\subsection{Limitations}

Several limitations qualify the findings of this work and delimit the conclusions properly drawn from them.

Each configuration was trained once, with a single seed, and none was repeated. The differences between closely ranked models --- in particular those separating the MAFF-Net family and the four ablated configurations --- are small relative to the per-image dispersion of the measures, so it is not possible, on the present evidence, to distinguish them from the variation between repeated runs. No statistical analysis was performed: no confidence intervals were estimated and no hypothesis tests conducted, and no claim of statistical significance is made anywhere in this work. Repeated runs across several seeds, with an appropriate interval estimate, would be required before the smaller differences reported here could be regarded as established.

The supervision strategy of Section III-H was active in every configuration, and its constituent terms were not ablated; the contribution of the boundary-focused supervision terms is therefore unmeasured, and no claim is advanced regarding their individual efficacy. The component ablation of Section V-C addresses the three architectural modules alone, on a single dataset, so the consistency of its findings across the other two benchmarks is not established.

The evaluation is confined to three dermoscopic benchmarks from the same challenge series. All images derive from a single imaging modality and annotation protocol, and the generality of the findings to other modalities, acquisition devices, or patient populations has not been assessed. The comparison is further restricted to models trained and evaluated within this study; no comparison is drawn against the literature, and several methods from Section II were not implemented here.

Finally, the relative standing within the MAFF-Net family should be stated plainly. The Restored MAFF-Net records the higher value on the majority of the measures in Sections V-A and V-B, including the Dice coefficient and the two distance measures on all three datasets. The evidence assembled here therefore does not establish the proposed configuration as the strongest member of the family on these benchmarks; consistent with the contributions stated in Section I, the value of this work lies in the combined region--boundary evaluation and the ranking disagreement it reveals, in the comparison with the standard baselines, and in the controlled analysis of the three components, not in a claim that the proposed configuration is the best available model. The qualitative observations of Section V-D rest on a few validation-split examples and are illustrative only.
